\section{Current Authorities: Who Teaches, Legislates, Interprets, and Judges?}

\subsection{Supreme authority}

The Roman Pontiff possesses supreme, full, immediate, and universal ordinary power in the Church and can always exercise it freely (CIC c.~331; CCEO c.~43). The college of bishops is also the subject of supreme and full power, but only with its head and never without him (CIC c.~336; CCEO c.~49). Ecumenical conciliar decrees require the approval, confirmation, and promulgation defined by law (CIC cc.~337--341; \sourceurl{https://www.vatican.va/archive/cod-iuris-canonici/eng/documents/cic_lib2-cann330-367_en.html}{CIC 330--367}).

Supreme power remains bounded by the divine constitution and law of the Church. ``Supreme'' does not mean that the Pope can make a contradiction true, dissolve what divine law makes indissoluble, or convert an intrinsically unjust act into justice. It means there is no higher human ecclesiastical authority or appeal from a papal sentence or decree (CIC c.~333 \S3).

\subsection{Particular Churches and intermediate bodies}

A diocesan bishop has ordinary, proper, and immediate power needed for his pastoral office except matters reserved by law. He governs with legislative, executive, and judicial power; he exercises legislative power personally and judicial power personally or through his tribunal (CIC cc.~381, 391; \sourceurl{https://www.vatican.va/archive/cod-iuris-canonici/eng/documents/cic_lib2-cann368-430_en.html}{CIC 368--430}). An eparchial bishop exercises corresponding authority within the Eastern canonical order.

Particular councils possess legislative power for their territory within universal law; their decrees require Apostolic See review before promulgation (CIC cc.~445--446). A Latin episcopal conference may issue general decrees only where universal law authorizes it or a special Apostolic See mandate grants it, with the voting, review, and promulgation conditions of c.~455 (\sourceurl{https://www.vatican.va/archive/cod-iuris-canonici/eng/documents/cic_lib2-cann431-459_en.html}{CIC 431--459}). Statements outside c.~455 do not thereby become legislation; their possible doctrinal authority is governed separately by CIC c.~753 and \sourceurl{https://www.vatican.va/content/john-paul-ii/en/motu_proprio/documents/hf_jp-ii_motu-proprio_22071998_apostolos-suos.html}{\emph{Apostolos suos}}, whose concluding norms require unanimity or a two-thirds plenary vote plus Apostolic See \emph{recognitio} for authentic conference magisterium and exclude commissions from exercising it.

Eastern patriarchal and major-archiepiscopal synods possess genuine legislative and judicial authority for their Churches \emph{sui iuris} (CCEO cc.~110, 152). They are not merely Eastern versions of Latin episcopal conferences. Synodal, patriarchal, eparchial, and common law must be classified on their own terms.

\subsection{The Roman Curia assists; it is not an autonomous legislator}

\emph{Praedicate Evangelium}, effective 5 June 2022, governs the Roman Curia (\sourceurl{https://www.vatican.va/content/francesco/en/apost_constitutions/documents/20220319-costituzione-ap-praedicate-evangelium.html}{text}). Curial institutions assist the Roman Pontiff in his mission. Article 30 provides that a curial institution cannot issue laws or general decrees having the force of law, or derogate from universal law, except in particular cases with a specific papal grant and approval. Executive decrees and instructions remain subordinate to law.

Leo XIV's General Regulations of the Roman Curia, effective 1 January 2026 \emph{ad experimentum} for five years, implement this structure and regulate drafting, review, papal approval, executive acts, and recourse (\sourceurl{https://www.vatican.va/content/leo-xiv/it/motu_proprio/documents/20251123-regolamento-generale-curia-romana.html}{official Italian}). Current institutional analysis must therefore distinguish a dicastery's competence, document genre, form of papal approval, promulgation, and intended juridical effect.

\subsection{Interpretation and doctrinal competence}

The Dicastery for Legislative Texts promotes knowledge and correct application of Latin and Eastern canon law, formulates authentic interpretations for specific papal approval, assists authorities, reviews certain normative texts, and studies legal gaps (\emph{Praedicate Evangelium} arts.~175--182; \sourceurl{https://www.vatican.va/content/romancuria/en/dicasteri/dicastero-testi-legislativi/profilo.html}{official profile}). Its formal \sourceurl{https://www.delegumtextibus.va/content/testilegislativi/it/interpretazioni-autentiche.html}{authentic-interpretation register} must be distinguished from explanatory notes or presidential responses.

The Dicastery for the Doctrine of the Faith safeguards doctrine concerning faith and morals and judges delicts reserved to it under proper norms (\emph{Praedicate Evangelium} arts.~69--78; \emph{Fidem servare}; \sourceurl{https://www.vatican.va/content/francesco/en/motu_proprio/documents/20220211-motu-proprio-fidem-servare.html}{text}). A doctrinal declaration can authoritatively identify a moral or divine-law boundary. It does not follow that every proposition in it is a new penal canon or that every form of papal approval converts it into legislation. Genre and promulgation still control juridical effect.

\subsection{Tribunals and administrative justice}

At first instance, the diocesan bishop is ordinarily the judge for cases not excepted by law, acting personally or through the judicial vicar and judges (CIC c.~1419). The Roman Rota ordinarily serves as a higher appellate tribunal at the Apostolic See and fosters unity of jurisprudence. The Apostolic Signatura is the supreme tribunal, supervises the administration of justice, and adjudicates specified recourses alleging violation of law by individual administrative acts (\emph{Praedicate Evangelium} arts.~189--204).

The Signatura's 2008 \sourceurl{https://www.vatican.va/content/benedict-xvi/en/apost_letters/documents/hf_ben-xvi_apl_20080621_antiqua-ordinatione.html}{\emph{Lex propria}, promulgated by \emph{Antiqua ordinatione}}, applies as amended by \sourceurl{https://www.vatican.va/content/francesco/en/motu_proprio/documents/20240228-motu-proprio-munus-tribunalis.html}{\emph{Munus Tribunalis}} in 2024. Judicial decisions can clarify law and carry persuasive or institutional weight, but a judgment is not thereby universal legislation. CIC c.~16 \S3 ordinarily limits a judicial or administrative interpretation to the persons and matter decided; CIC c.~19 gives jurisprudence a role in filling nonpenal gaps.

Administrative recourse likewise distinguishes the merits of a pastoral dispute from whether an authority violated law or procedure. Written reasons, hearing affected persons, competence, notification, deadlines, suspension, hierarchical recourse, and Signatura review are positive-law means of serving justice. Invoking a divine purpose does not excuse procedural illegality.

\subsection{A norm's institutional path}

\begin{center}
\begin{tikzpicture}[
  node distance=6mm and 7mm,
  box/.style={draw, rounded corners=1pt, align=center, inner sep=5pt, text width=0.25\linewidth, font=\footnotesize},
  arrow/.style={-{Latex[length=2mm]}, line width=0.45pt}
]
\node[box] (source) {Identify the source\\divine claim, lawgiver, custom, treaty, or received civil norm};
\node[box, right=of source] (form) {Verify juridical form\\competence, promulgation, scope, effective date, amendments};
\node[box, right=of form] (meaning) {Interpret\\text, context, purpose, authentic responses, canonical tradition};
\node[box, below=of meaning] (apply) {Apply or adjudicate\\facts, rights, proof, procedure, equity, remedies};
\node[box, left=of apply] (review) {Review\\administrative recourse, appeal, jurisprudence, legislative correction};
\node[box, left=of review] (audit) {Re-audit\\higher and divine law, civil interaction, changed facts or law};
\draw[arrow] (source) -- (form);
\draw[arrow] (form) -- (meaning);
\draw[arrow] (meaning) -- (apply);
\draw[arrow] (apply) -- (review);
\draw[arrow] (review) -- (audit);
\draw[arrow] (audit) -- (source);
\end{tikzpicture}
\end{center}

The diagram is deliberately circular. Application exposes ambiguity or injustice; recourse and jurisprudence may reveal a gap; legislation can change; and every new act requires renewed source classification. No office substitutes for the audit.
